\chapter{Introduction}
\begin{figure}[h]
  \centering
  \includegraphics[width=\textwidth]{./images/logos/he-arc-banner.jpg}
  \caption{Teaser figure.}
  \label{Fig_teaser}
\end{figure}

\section{Comment formater le texte}
\textbf{Comment mettre un texte en gras.}
\underline{Comment souligner un texte.} 
\textit{Comment mettre un texte en italique.}
\underline{\textbf{\textit{On peut aussi mélanger les 3.}}}

Pour commencer un nouveau paragraphe, il suffit d'ajouter 1 ligne vide le fichier .tex.

Pour faire une liste :
\begin{itemize}
    \item élément 1
    \item élément 2
    \item élément 3
\end{itemize}

Pour faire une liste numérotée :
\begin{enumerate}
    \item élément 1
    \item élément 2
    \item élément 3
\end{enumerate}

Pour rajouter une figure et la référencer (voir Figure \ref{Fig_exemple}).
\begin{figure}[h]
    \centering
    \includegraphics[width=0.5\textwidth]{./images/logos/he-arc-logo.png}
    \caption{Ceci est le nouveau logo de la HE-Arc.}
    \label{Fig_exemple}
\end{figure}

\section{Lorem ipsum}
\lipsum[1-1] % PLACEHOLDER
\subsection{Dolor sit}
\lipsum[1-1] % PLACEHOLDER
\subsubsection{amet}
\lipsum[1-1] % PLACEHOLDER
\paragraph{consectetuer}
\lipsum[1-1] % PLACEHOLDER